% What survives the pushforward at each resolution (GBDT, 100 pilot files, interventional TreeSHAP as the reference;
% byte deletion, runs 66 and 74). Data: scripts/paper_figures.py -> data/pf_gbdt.dat: per source resolution, the median over
% files of the pushed-forward ranking's deletion area over that of TreeSHAP's pullback. "bytes" is drawn at 2^20 windows.
\begin{figure}[t]
  \centering
  \begin{tikzpicture}
    \begin{axis}[
      paper,
      xmode=log, log basis x=2, xmin=40, xmax=1700000, ymin=0.5, ymax=1.05,
      xtick={64,256,1024,4096,16384,65536,262144,1048576},
      xticklabels={64,256,1k,4k,16k,64k,256k,bytes},
      xlabel={Resolution of the source explanation (windows)},
      ylabel={Evidence removed (relative to TreeSHAP)},
      legend pos=south east, legend style={font=\scriptsize},
    ]
      \addplot[ours, thick, mark=*, mark size=1.6pt] table[x=windows, y=trend] {assets/data/pf_gbdt.dat};
      \addlegendentry{TreeSHAP's information, pushed forward}
      \addplot[direct!60!black, thick, mark=triangle*, mark size=1.8pt] table[x=windows, y=kernelshap] {assets/data/pf_gbdt.dat};
      \addlegendentry{KernelSHAP, pushed forward}
      \addplot[direct!60!black, thick, dashed, mark=diamond*, mark size=1.8pt] table[x=windows, y=adaptive] {assets/data/pf_gbdt.dat};
      \addlegendentry{Adaptive KernelSHAP, pushed forward}
      \addplot[direct, thick, mark=square*, mark size=1.4pt] table[x=windows, y=occlusion] {assets/data/pf_gbdt.dat};
      \addlegendentry{Occlusion, pushed forward}
    \end{axis}
  \end{tikzpicture}
  \caption{Faithfulness of the pushforward rises with the resolution of the explanation it transfers (GBDT, 100 files).
  The evidence removed by deleting the bytes of the top-ranked features, relative to deleting the bytes that TreeSHAP's
  pullback ranks highest (median over files). TreeSHAP's own explanation, expressed over windows of each size and pushed
  forward, shows what the information at each resolution supports; KernelSHAP and adaptive KernelSHAP follow it or
  exceed it, and occlusion, which changes the features too little in small windows, stays flat.}
  \label{fig:XML:PushforwardResolution}
\end{figure}
